\begin{helper}
Let $Q$ be a three-element finite set of vertices, and let $F$ be any
real-valued function on finite vertex sets.  Then there are distinct
vertices $a,b,c\in Q$ such that $Q=\{a,b,c\}$, the two-element subsets of
$Q$ are exactly
\[
\{a,b\},\qquad \{a,c\},\qquad \{b,c\},
\]
and
\[
\sum_P \mathbf{1}_{\{|P|=2,\ P\subseteq Q\}} F(P)
=F(\{a,b\})+F(\{a,c\})+F(\{b,c\}).
\]
\end{helper}
\begin{proof}
Since $|Q|=3$, choose an enumeration $Q=\{a,b,c\}$ with
$a,b,c$ pairwise distinct.  Plainly $a,b,c\in Q$.

Let $P\subseteq Q$ have size two.  Write $P=\{x,y\}$ with $x\ne y$.
Because $x,y\in Q=\{a,b,c\}$, each of $x$ and $y$ is one of $a,b,c$.
The case $x=y$ is impossible, and the remaining ordered possibilities give
exactly the unordered sets $\{a,b\}$, $\{a,c\}$, and $\{b,c\}$.  Conversely,
each of these three displayed sets has size two, by pairwise distinctness,
and is contained in $Q$.  Hence they are precisely the complete collection
of two-element subsets of $Q$.

The indicated sum has a nonzero summand exactly on this complete collection
of three subsets and is zero on every other finite set $P$.  Therefore it
reduces to the three displayed terms
$F(\{a,b\})+F(\{a,c\})+F(\{b,c\})$.
\end{proof}
